% Working outline agreed 29 October 2026. Target: four pages including references.
\documentclass[11pt,twoside]{article}
\usepackage{asp2023}
\aspSuppressVolSlug
\resetcounters
\bibliographystyle{asp2021}
\markboth{Smallwood et al.}{Spatial Filtering in Practice}
\begin{document}
\title{Spatial Filtering in Practice: A GPU Pipeline for Real-Time RFI Mitigation}
% TODO: confirm the complete author list, affiliations, email, ADS metadata and index entries.
\author{Justin (Jay) Smallwood et al.}
\affil{Author details to be confirmed.}

\begin{abstract}
\textit{Writing outline:} Introduce the opportunity for GPU software to make
spatial filtering practical to develop and integrate into telescope systems.
Present Starspin, with short-time correlation and eigendecomposition built
into its signal chain. Summarise the measured processing cost and one
LAMBDA example, distinguishing live filtering from offline protection results.
Conclude with the scope for richer methods and prior information.
% Replace this prompt with the final abstract after fresh benchmarks.
% ADASS-abstract.tex records the submission; its throughput claims need rechecking.
\end{abstract}

\section{Motivation: Making Spatial Filtering Practical}
\begin{itemize}
\item Introduce radio-frequency interference (RFI), the cost of discarding
contaminated data, and the aim of retaining astronomy while suppressing RFI.
\item Explain the opportunity: GPU compute, including tensor-core correlation
and beamforming, provides a real-time processing engine with software
flexibility and faster feedback between ideas, implementation and evaluation.
\item Position Starspin as a pipeline designed with RFI mitigation as a core
requirement. Briefly acknowledge earlier spatial-filtering implementations;
avoid claims of first use or an FPGA--GPU contest.
\end{itemize}

\section{From Array Signals to Spatial Filtering}
\begin{itemize}
\item For a software audience, introduce antenna voltages, correlations
(visibilities), and spatial modes before using them.
\item Explain eigendecomposition and simple mode nulling. Eigenmodes can mix
astronomy and interference; a strong mode is not automatically RFI.
\item Use one LAMBDA example to show both the opportunity and the need to
preserve astronomical signals. Cite the submitted protection paper for detail.
\end{itemize}
\noindent\textit{Figure placeholder:} Compact real-data comparison showing
interference, filtering and astronomical protection. Identify the offline
analysis explicitly; it is separate from the live pipeline measurement.

\section{Starspin: Mitigation in the Signal Chain}
\begin{itemize}
\item Trace input capture, unpacking/alignment, fine channelisation and
correction, followed by correlation and beamforming.
\item Show short-time correlations feeding eigendecomposition and adaptive
beam-weight updates. Explain why the update interval must follow changes
in the interferer; do not assume a universal cadence.
\item Explain the voltage path and weight-update path separately. Identify
what is implemented and where richer filtering methods can enter.
\item Describe tensor-core libraries and GPU design choices only as needed
to explain processing capacity and the ability to iterate on methods.
\end{itemize}
\noindent\textit{Figure placeholder:} Signal-chain diagram with the
correlation--eigendecomposition--weight-update loop and beam/visibility outputs.

\section{Performance and the Resolution of Filter Estimation}
\begin{itemize}
\item Rerun sustained-throughput benchmarks after the fine-channelisation
changes. Compare eigendecomposition off/on with other settings held fixed;
state whether the enabled path includes weight updates and filter application.
\item Record GPU, input count and polarisation, channel widths, covariance
integration and update intervals, enabled stages, run duration and data loss.
Separate processed bandwidth, input byte rate and update latency.
\item Discuss the open choice of estimating filters per fine channel, per
coarse channel, or through a mixed approach. Distinguish estimation resolution
from the resolution at which filtering is applied.
\item Evaluate the trade-off among throughput, frequency-dependent spatial
structure and the cadence needed for moving interferers. A faster estimate
must still provide useful filtering and retain astronomy.
\end{itemize}
\noindent\textit{Table/plot placeholder:} Fresh paired throughput results,
with configuration and timing context. Coarse/fine/mixed comparisons remain
planned until implemented and measured. Do not reuse old numbers as current.

\section{Richer Filters and Further Work}
\begin{itemize}
\item Extend the progression from nulling to eigenvalue shrinking and
astronomical protection. State what each changes and why it is useful;
refer to the submitted paper for the full protection methods.
\item Introduce satellite two-line elements (TLEs) as a future source of
prior information to guide interference estimation or selection. Geometry
alone does not establish a calibrated antenna-space signature.
\item Close with the opportunity to investigate better filtering in a
working telescope environment. Keep demonstrated capabilities, offline
results and proposed developments distinct.
\end{itemize}

% TODO: add a small set of verified references to C802.bib, then enable:
% \bibliography{C802}
% The supplied C802.bib still contains template examples, not selected references.
% TODO: acknowledgements and subject/software index entries before submission.
% \bookpartphoto[width=1.0\textwidth]{foobar.eps}{FooBar Photo (Photo: Any Photographer)}
\end{document}
